\documentclass[10pt,a4paper]{amsart}
\usepackage{amsmath,amssymb,amsthm,graphicx}
\usepackage{verbatim}
\newenvironment{reference}{\comment}{\endcomment}
%\newenvironment{reference}{}{}
\newenvironment{slogan}{\comment}{\endcomment}
\newenvironment{history}{\comment}{\endcomment}

% For commutative diagrams we use Xy-pic
\usepackage[all]{xy}

% We use 2cell for 2-commutative diagrams.
\xyoption{2cell}
% Theorem environments.
%
\theoremstyle{plain}
\newtheorem{theorem}[subsection]{Theorem}
\newtheorem{proposition}[subsection]{Proposition}
\newtheorem{lemma}[subsection]{Lemma}

\theoremstyle{definition}
\newtheorem{definition}[subsection]{Definition}
\newtheorem{example}[subsection]{Example}
\newtheorem{exercise}[subsection]{Exercise}
\newtheorem{situation}[subsection]{Situation}

\theoremstyle{remark}
\newtheorem{remark}[subsection]{Remark}
\newtheorem{remarks}[subsection]{Remarks}

\numberwithin{equation}{subsection}

% Macros
%
\def\lim{\mathop{\mathrm{lim}}\nolimits}
\def\colim{\mathop{\mathrm{colim}}\nolimits}
\def\Spec{\mathop{\mathrm{Spec}}}
\def\Hom{\mathop{\mathrm{Hom}}\nolimits}
\def\Ext{\mathop{\mathrm{Ext}}\nolimits}
\def\SheafHom{\mathop{\mathcal{H}\!\mathit{om}}\nolimits}
\def\SheafExt{\mathop{\mathcal{E}\!\mathit{xt}}\nolimits}
\def\Sch{\mathit{Sch}}
\def\Mor{\mathop{\mathrm{Mor}}\nolimits}
\def\Ob{\mathop{\mathrm{Ob}}\nolimits}
\def\Sh{\mathop{\mathit{Sh}}\nolimits}
\def\NL{\mathop{N\!L}\nolimits}
\def\CH{\mathop{\mathrm{CH}}\nolimits}
\def\proetale{{pro\text{-}\acute{e}tale}}
\def\etale{{\acute{e}tale}}
\def\QCoh{\mathit{QCoh}}
\def\Ker{\mathop{\mathrm{Ker}}}
\def\Im{\mathop{\mathrm{Im}}}
\def\Coker{\mathop{\mathrm{Coker}}}
\def\Coim{\mathop{\mathrm{Coim}}}

% Boxtimes
%
\DeclareMathSymbol{\boxtimes}{\mathbin}{AMSa}{"02}

%
% Macros for moduli stacks/spaces
%
\def\QCohstack{\mathcal{QC}\!\mathit{oh}}
\def\Cohstack{\mathcal{C}\!\mathit{oh}}
\def\Spacesstack{\mathcal{S}\!\mathit{paces}}
\def\Quotfunctor{\mathrm{Quot}}
\def\Hilbfunctor{\mathrm{Hilb}}
\def\Curvesstack{\mathcal{C}\!\mathit{urves}}
\def\Polarizedstack{\mathcal{P}\!\mathit{olarized}}
\def\Complexesstack{\mathcal{C}\!\mathit{omplexes}}
% \Pic is the operator that assigns to X its picard group, usage \Pic(X)
% \Picardstack_{X/B} denotes the Picard stack of X over B
% \Picardfunctor_{X/B} denotes the Picard functor of X over B
\def\Pic{\mathop{\mathrm{Pic}}\nolimits}
\def\Picardstack{\mathcal{P}\!\mathit{ic}}
\def\Picardfunctor{\mathrm{Pic}}
\def\Deformationcategory{\mathcal{D}\!\mathit{ef}}
\usepackage{cleveref}
\crefname{originalreference}{original reference}{original references}
\Crefname{originalreference}{Original reference}{Original references}
\makeatletter
\newlabel{lemma-nr-roots-unchanged}{{09H4}{source}{}{sourcepointer-1}{}}
\newlabel{lemma-nr-roots-unchanged@cref}{{[originalreference][0][]09H4}{[1][0][]source}{}{}{}}
\makeatother
\crefname{originalreference}{original reference}{original references}
\Crefname{originalreference}{Original reference}{Original references}
\makeatletter
\newlabel{lemma-take-pth-root}{{09HF}{source}{}{sourcepointer-2}{}}
\newlabel{lemma-take-pth-root@cref}{{[originalreference][0][]09HF}{[1][0][]source}{}{}{}}
\makeatother
\title{A finite purely inseparable extension is an iterated tower of p-th root adjunctions of degree p}
\author{The Stacks Project Authors}
\date{stacks-project:fields.tex:local-snapshot-20261012; GNU FDL 1.2 or later}
\begin{document}
\maketitle
\noindent\textit{Editorial note.} All delivered mathematical text is the author's own, copied from the pinned Stacks Project chapter file fields.tex; no statement or proof step was rewritten, reordered or removed.\par\medskip
\noindent\textit{Editorial note.} The counted proof is the author's complete argument: The tower produced here is the input for the purely inseparable case of lemma-separable-degree (delivered L48344).\par\medskip
\noindent\textit{Editorial note.} The label lemma-finite-purely-inseparable names the counted statement, and the proof environment immediately following it in the source is delivered verbatim as its proof.\par\medskip
\noindent\textit{Editorial note.} The counted statement is carried by the author's own lemma environment in the chapter source.\par\medskip
\noindent\textit{Allowed standard prerequisites (editorial).}\begin{itemize}
\item Same chapter: lemma-nr-roots-unchanged (the number of p-th roots of an element is unchanged under field extension - used to see t = alpha\textasciicircum{}p has no p-th root in F), lemma-take-pth-root (adjoining a p-th root of an element with no p-th root gives a degree p purely inseparable extension); definition of purely inseparable (every element has a p-power power in the base)
\end{itemize}\medskip
\ifdefined\hypertarget\hypertarget{sourcepointer-1}{}\fi
\noindent\textit{Source reference.} Original reference 09H4: sources/stacks-fields/pinned-original-source.tex, source lines 1442--1442. Prerequisite pointer invoked at 1911; it names a standard fact of the same chapter that the counted argument uses, and every step of the counted proof built on it is retained verbatim in the delivered body.\par\medskip
\ifdefined\hypertarget\hypertarget{sourcepointer-2}{}\fi
\noindent\textit{Source reference.} Original reference 09HF: sources/stacks-fields/pinned-original-source.tex, source lines 1834--1834. Prerequisite pointer invoked at 1913; it names a standard fact of the same chapter that the counted argument uses, and every step of the counted proof built on it is retained verbatim in the delivered body.\par\medskip
\begin{lemma}
\label{lemma-finite-purely-inseparable}
Let $E/F$ be a finite purely inseparable field extension of
characteristic $p > 0$. Then there exists a sequence of elements
$\alpha_1, \ldots, \alpha_n \in E$ such that we obtain a tower
of fields
$$
E = F(\alpha_1, \ldots, \alpha_n) \supset
F(\alpha_1, \ldots, \alpha_{n - 1}) \supset
\ldots
\supset F(\alpha_1) \supset F
$$
such that each intermediate extension is of degree $p$ and comes
from adjoining a $p$th root. Namely,
$\alpha_i^p \in F(\alpha_1, \ldots, \alpha_{i - 1})$
is an element which does not have a $p$th root in
$F(\alpha_1, \ldots, \alpha_{i - 1})$ for $i = 1, \ldots, n$.
\end{lemma}

\begin{proof}
By induction on the degree of $E/F$. If the degree of the extension is $1$
then the result is clear (with $n = 0$). If not, then choose
$\alpha \in E$, $\alpha \not \in F$. Say $\alpha^{p^r} \in F$ for some
$r > 0$. Pick $r$ minimal and replace $\alpha$ by $\alpha^{p^{r - 1}}$.
Then $\alpha \not \in F$, but $\alpha^p \in F$. Then $t = \alpha^p$ is not
a $p$th power in $F$ (because that would imply $\alpha \in F$, see
Lemma \ref{lemma-nr-roots-unchanged} or its proof).
Thus $F \subset F(\alpha)$ is a subextension of degree $p$
(Lemma \ref{lemma-take-pth-root}). By induction we find
$\alpha_1, \ldots, \alpha_n \in E$ generating $E/F(\alpha)$
satisfying the conclusions of the lemma.
The sequence $\alpha, \alpha_1, \ldots, \alpha_n$ does the job
for the extension $E/F$.
\end{proof}

\end{document}
